\documentclass{amsart}
\usepackage{amsmath,amssymb,amsthm}

\newtheorem{theorem}{Theorem}
\newtheorem{corollary}[theorem]{Corollary}
\newtheorem{lemma}[theorem]{Lemma}

\newcommand{\assign}{\coloneqq}
\newcommand{\abs}[1]{\left|#1\right|}
\newcommand{\norm}[1]{\left\|#1\right\|}

\title{Convergence of the Autocovariance of the Normalized Zeta-Process}
\author{Stephen Crowley}
\date{April 2026}

\begin{document}
\maketitle

\section{Notation and Setup}

Let $\zeta(s)$ denote the Riemann zeta function and let
\[
\vartheta(t) \assign \operatorname{Im}\log\Gamma\!\left(\tfrac{1}{4}+\tfrac{it}{2}\right) - \tfrac{t}{2}\log\pi
\]
be the Riemann--Siegel theta function. The function $\vartheta$ is real-analytic
on $(0,\infty)$, attains its global minimum on $(0,\infty)$ at a unique point
$t_0$, and satisfies $\vartheta'(t) > 0$ for all $t > t_0$. Set $\tau_0 \assign
\vartheta(t_0)$. Let $\sigma : [\tau_0,\infty) \to [t_0,\infty)$ denote the
inverse of $\vartheta\big|_{[t_0,\infty)}$, so $\vartheta(\sigma(\tau))=\tau$
and $\sigma'(\tau) = 1/\vartheta'(\sigma(\tau))$ for all $\tau \ge \tau_0$.

The Stirling expansion gives
\begin{equation}\label{eq:theta-prime}
  \vartheta'(t) = \tfrac{1}{2}\log\!\left(\tfrac{t}{2\pi}\right) + O(t^{-1}),
  \qquad \vartheta''(t) = O(t^{-1}), \qquad \vartheta'''(t) = O(t^{-2}).
\end{equation}

Define the second-moment integral
\[
  V(T) \assign \int_0^T \abs{\zeta\!\left(\tfrac{1}{2}+it\right)}^2 dt.
\]
By the classical result of Ingham \cite{Ingham1928},
\begin{equation}\label{eq:Ingham}
  V(T) = T\log T - T(1 + \log 2\pi) + O(T^{1/2+\varepsilon})
  \quad\text{for every } \varepsilon > 0.
\end{equation}

From $\vartheta(t) \sim \tfrac{t}{2}\log(t/2\pi)$ one obtains
$\sigma(\tau) \sim 2\tau/\log(2\tau/2\pi) \sim 2\tau/\log\tau$, and then
\begin{equation}\label{eq:sigma-asymp}
  \sigma(\tau) = \frac{2\tau}{\log\tau}\left(1 + O\!\left(\frac{\log\log\tau}{\log\tau}\right)\right).
\end{equation}

Define the unitary pullback
\[
  Y(\tau) \assign \frac{\zeta\!\left(\tfrac{1}{2}+i\sigma(\tau)\right)}{\sqrt{\vartheta'(\sigma(\tau))}},
\]
the cumulative local energy
\[
  W(\tau) \assign \int_{\tau_0}^{\tau} \abs{Y(s)}^2\,ds,
\]
and the normalized process
\[
  X(\tau) \assign Y(\tau)\sqrt{\frac{\tau}{W(\tau)}}.
\]

\begin{lemma}[Energy identity]\label{lem:energy}
For all $\tau \ge \tau_0$,
\[
  W(\tau) = V(\sigma(\tau)) - V(\sigma(\tau_0)).
\]
\end{lemma}
\begin{proof}
Substitute $s = \vartheta(t)$, so $t = \sigma(s)$ and $ds = \vartheta'(t)\,dt$:
\[
  W(\tau) = \int_{\tau_0}^{\tau} \frac{\abs{\zeta(\tfrac{1}{2}+i\sigma(s))}^2}{\vartheta'(\sigma(s))}\,ds
  = \int_{\sigma(\tau_0)}^{\sigma(\tau)} \abs{\zeta(\tfrac{1}{2}+it)}^2\,dt
  = V(\sigma(\tau)) - V(\sigma(\tau_0)). \qedhere
\]
\end{proof}

\begin{lemma}[Energy asymptotics]\label{lem:W-asymp}
\[
  W(\tau) = 2\tau\!\left(1 + O\!\left(\frac{\log\log\tau}{\log\tau}\right)\right).
\]
\end{lemma}
\begin{proof}
From \eqref{eq:Ingham} and \eqref{eq:sigma-asymp},
\[
  V(\sigma(\tau)) \sim \sigma(\tau)\log\sigma(\tau)
  \sim \frac{2\tau}{\log\tau}\cdot\log\!\left(\frac{2\tau}{\log\tau}\right)
  = \frac{2\tau}{\log\tau}\cdot\left(\log\tau + O(\log\log\tau)\right)
  = 2\tau\!\left(1 + O\!\left(\frac{\log\log\tau}{\log\tau}\right)\right).
\]
The term $V(\sigma(\tau_0))$ is a fixed constant. The error term in
\eqref{eq:Ingham} contributes $O(\sigma(\tau)^{1/2+\varepsilon}) =
O(\tau^{1/2+\varepsilon}/(\log\tau)^{1/2+\varepsilon}) = o(\tau)$, which is
absorbed into the stated error.
\end{proof}

Write $\varepsilon(\tau) \assign W(\tau)/(2\tau) - 1 = O(\log\log\tau/\log\tau)$.

\section{Main Theorem}

\begin{theorem}\label{thm:cov}
For every fixed $h \in \mathbb{R}$, the Ces\`aro autocovariance
\[
  R(h) \assign \lim_{T\to\infty} \frac{1}{T}\int_{\tau_0}^T
  X(\tau+h)\,\overline{X(\tau)}\,d\tau
\]
exists and equals
\[
  R(h) = \frac{e^{-ih}\sin h}{h},
\]
with the convention $R(0) = 1$.
\end{theorem}

\begin{corollary}\label{cor:spectral}
The spectral density of $X$ is $\hat{S}(\omega) = \pi\cdot\mathbf{1}_{[-2,0]}(\omega)$.
\end{corollary}
\begin{proof}
Using the standard identity $\int_{-\infty}^{\infty}(\sin h/h)\,e^{-i\alpha h}\,dh
= \pi\cdot\mathbf{1}_{[-1,1]}(\alpha)$:
\[
  \hat{S}(\omega)
  = \int_{-\infty}^{\infty} \frac{e^{-ih}\sin h}{h}\,e^{-i\omega h}\,dh
  = \int_{-\infty}^{\infty} \frac{\sin h}{h}\,e^{-i(\omega+1)h}\,dh
  = \pi\cdot\mathbf{1}_{[-1,1]}(\omega+1)
  = \pi\cdot\mathbf{1}_{[-2,0]}(\omega). \qedhere
\]
\end{proof}

\section{Proof of Theorem~\ref{thm:cov}}

\subsection*{Step 1: Reduction to $\mathcal{I}(T,h)$}

Substituting the definitions,
\[
  X(\tau+h)\,\overline{X(\tau)}
  = Y(\tau+h)\,\overline{Y(\tau)}\cdot
    \frac{\sqrt{(\tau+h)\tau}}{\sqrt{W(\tau+h)W(\tau)}}.
\]
Define
\[
  \alpha(\tau,h) \assign \frac{\sqrt{(\tau+h)\tau}}{\sqrt{W(\tau+h)W(\tau)}}.
\]

\begin{lemma}\label{lem:alpha}
For each fixed $h \in \mathbb{R}$, $\alpha(\tau,h) = \tfrac{1}{2} + O(\log\log\tau/\log\tau)$ as $\tau\to\infty$.
\end{lemma}
\begin{proof}
By Lemma~\ref{lem:W-asymp},
\[
  \alpha(\tau,h)
  = \frac{\sqrt{(\tau+h)\tau}}{\sqrt{2(\tau+h)(1+\varepsilon(\tau+h))\cdot 2\tau(1+\varepsilon(\tau))}}
  = \frac{1}{2}\cdot\frac{1}{\sqrt{(1+\varepsilon(\tau+h))(1+\varepsilon(\tau))}}.
\]
Since both $\varepsilon(\tau)$ and $\varepsilon(\tau+h) = O(\log\log\tau/\log\tau)$, the
expansion $(1+u)^{-1/2} = 1 - u/2 + O(u^2)$ gives the result.
\end{proof}

Write $\alpha = \tfrac{1}{2} + \beta(\tau,h)$ with $\beta = O(\log\log\tau/\log\tau)$, so
\begin{equation}\label{eq:split}
  \frac{1}{T}\int_{\tau_0}^T X(\tau+h)\,\overline{X(\tau)}\,d\tau
  = \tfrac{1}{2}\,\mathcal{I}(T,h) + \mathcal{E}(T,h),
\end{equation}
where
\[
  \mathcal{I}(T,h) \assign \frac{1}{T}\int_{\tau_0}^T Y(\tau+h)\,\overline{Y(\tau)}\,d\tau,
  \qquad
  \mathcal{E}(T,h) \assign \frac{1}{T}\int_{\tau_0}^T \beta(\tau,h)\,Y(\tau+h)\,\overline{Y(\tau)}\,d\tau.
\]

\begin{lemma}\label{lem:E-vanishes}
$\mathcal{E}(T,h) \to 0$ as $T\to\infty$ for every fixed $h$.
\end{lemma}
\begin{proof}
Let $B \assign \sup_{\tau \ge \tau_0}|\beta(\tau,h)| = O(\log\log\tau_0/\log\tau_0)$,
which is a fixed finite constant. By Cauchy--Schwarz,
\[
  |\mathcal{E}(T,h)|
  \le B \left(\frac{1}{T}\int_{\tau_0}^T |Y(\tau+h)|^2\,d\tau\right)^{1/2}
       \left(\frac{1}{T}\int_{\tau_0}^T |Y(\tau)|^2\,d\tau\right)^{1/2}.
\]
By Lemma~\ref{lem:W-asymp}, $T^{-1}\int_{\tau_0}^T |Y|^2\,d\tau = W(T)/T \to 2$.
The shifted factor satisfies $T^{-1}\int_{\tau_0}^T |Y(\tau+h)|^2\,d\tau
= (W(T+h)-W(h+\tau_0))/T \to 2$ likewise. Hence $|\mathcal{E}(T,h)| \le B \cdot O(1)$.
Since $B \to 0$ as $\tau_0\to\infty$ and is fixed once $\tau_0$ is chosen,
and the $O(1)$ factor is uniform in $T$, for any $\eta > 0$ we may choose
$\tau_0$ large enough that $B < \eta$, giving $\limsup_{T\to\infty}|\mathcal{E}(T,h)| < 2\eta$.
Since $\eta$ is arbitrary, $\mathcal{E}(T,h)\to 0$.
\end{proof}

By \eqref{eq:split} and Lemma~\ref{lem:E-vanishes}, it suffices to show
\begin{equation}\label{eq:suffices}
  \lim_{T\to\infty}\mathcal{I}(T,h) = \int_{-2}^0 e^{i\omega h}\,d\omega
  = \frac{2e^{-ih}\sin h}{h},
\end{equation}
since then $R(h) = \tfrac{1}{2}\cdot\tfrac{2e^{-ih}\sin h}{h} = e^{-ih}\sin h/h$.

\subsection*{Step 2: Change of Variables and Height Shift}

Set $t = \sigma(\tau)$, so $\tau = \vartheta(t)$ and $d\tau = \vartheta'(t)\,dt$.
Let $T_t \assign \sigma(T)$ denote the upper limit after the substitution.
For the shifted term, define
\begin{equation}\label{eq:delta-def}
  \delta_h(t) \assign \sigma(\vartheta(t)+h) - t.
\end{equation}

\begin{lemma}\label{lem:delta}
For each fixed $h$,
\[
  \delta_h(t) = \frac{h}{\vartheta'(t)} + O\!\left(\frac{1}{t(\log t)^3}\right)
\]
as $t\to\infty$.
\end{lemma}
\begin{proof}
Apply the Taylor expansion of $\sigma$ at $\vartheta(t)$:
\[
  \sigma(\vartheta(t)+h) - \sigma(\vartheta(t))
  = h\,\sigma'(\vartheta(t)) + \tfrac{h^2}{2}\sigma''(\xi)
\]
for $\xi$ between $\vartheta(t)$ and $\vartheta(t)+h$. Since $\sigma(\vartheta(t))=t$
and $\sigma'(\vartheta(t)) = 1/\vartheta'(t)$, the leading term is $h/\vartheta'(t)$.
Differentiating $\sigma'(s) = 1/\vartheta'(\sigma(s))$ gives
$\sigma''(s) = -\vartheta''(\sigma(s))/(\vartheta'(\sigma(s)))^3$.
By \eqref{eq:theta-prime}, $\vartheta''(t) = O(t^{-1})$ and $\vartheta'(t)
\sim \tfrac{1}{2}\log(t/2\pi)$, so $\sigma''(\xi) = O(1/(t(\log t)^3))$.
\end{proof}

In particular $\delta_h(t) \sim 2h/\log(t/2\pi) \to 0$. After the substitution $\tau\to t$,
\begin{equation}\label{eq:I-height}
  \mathcal{I}(T,h) = \frac{1}{\vartheta(T_t)}\int_{t_0}^{T_t}
  \frac{\zeta(\tfrac{1}{2}+i(t+\delta_h(t)))\,\overline{\zeta(\tfrac{1}{2}+it)}}
       {\sqrt{\vartheta'(t+\delta_h(t))\,\vartheta'(t)}}\,\vartheta'(t)\,dt.
\end{equation}

\begin{lemma}\label{lem:jacobian}
For fixed $h$,
\[
  \mathcal{I}(T,h)
  = \frac{1}{\vartheta(T_t)}\int_{t_0}^{T_t}
    \zeta\!\left(\tfrac{1}{2}+i(t+\delta_h(t))\right)
    \overline{\zeta\!\left(\tfrac{1}{2}+it\right)}\,dt
  \;+\; o(1)
\]
as $T_t\to\infty$.
\end{lemma}
\begin{proof}
From \eqref{eq:I-height}, after cancelling one factor of $\sqrt{\vartheta'(t)}$,
the integrand carries the ratio $\sqrt{\vartheta'(t)}/\sqrt{\vartheta'(t+\delta_h(t))}
= 1 + \gamma(t)$ where, by \eqref{eq:theta-prime} and Lemma~\ref{lem:delta},
\[
  \gamma(t) = \frac{\vartheta'(t+\delta_h(t))-\vartheta'(t)}{2\vartheta'(t)}
             + O(\gamma(t)^2)
  = \frac{\delta_h(t)\,\vartheta''(t)}{2\vartheta'(t)} + O\!\left(\frac{1}{(t\log t)^2}\right)
  = O\!\left(\frac{1}{t\log t}\right).
\]
Let $M \assign \sup_{t\ge t_0}|\gamma(t)|$, a finite constant. By Cauchy--Schwarz,
\[
  \frac{1}{\vartheta(T_t)}\int_{t_0}^{T_t} |\gamma(t)|\,
  |\zeta(\tfrac{1}{2}+i(t+\delta_h(t)))|\,|\zeta(\tfrac{1}{2}+it)|\,dt
  \le M\cdot\frac{V(T_t)^{1/2}\cdot V(T_t+|h|)^{1/2}}{\vartheta(T_t)}.
\]
By \eqref{eq:Ingham}, $V(T_t) \sim T_t\log T_t$ and $\vartheta(T_t) \sim
\tfrac{T_t}{2}\log T_t$, so the right side is $O(M)$. Since $M = O(1/(t_0\log t_0))$,
for any $\eta > 0$ choose $t_0$ large enough that $M < \eta$; then for all $T_t$
the error is at most $C\eta$. As $\eta$ is arbitrary the error is $o(1)$.
\end{proof}

\subsection*{Step 3: Approximate Functional Equation and Diagonal Extraction}

For $t \ge t_{\mathrm{RS}}$ the Riemann--Siegel formula \cite{Titchmarsh1986} gives
\begin{equation}\label{eq:RS}
  \zeta\!\left(\tfrac{1}{2}+it\right)
  = \underbrace{\sum_{n\le N(t)} n^{-1/2-it}}_{A(t)}
  + e^{-2i\vartheta(t)}\underbrace{\sum_{n\le N(t)} n^{-1/2+it}}_{B(t)}
  + R(t),
\end{equation}
where $N(t) = \lfloor\sqrt{t/(2\pi)}\rfloor$ and $|R^{(j)}(t)| \le C_j\,t^{-1/4-j/2}$
for all $j\ge 0$.

Write $\delta = \delta_h(t)$ throughout. The product
$\zeta(\tfrac{1}{2}+i(t+\delta))\,\overline{\zeta(\tfrac{1}{2}+it)}$ expands
via \eqref{eq:RS} into four main terms:
\begin{align*}
  \mathrm{I}   &= A(t+\delta)\,\overline{A(t)}, \\
  \mathrm{II}  &= A(t+\delta)\,e^{2i\vartheta(t)}\,\overline{B(t)}, \\
  \mathrm{III} &= e^{-2i\vartheta(t+\delta)}\,B(t+\delta)\,\overline{A(t)}, \\
  \mathrm{IV}  &= e^{-2i\vartheta(t+\delta)}\,B(t+\delta)\,e^{2i\vartheta(t)}\,\overline{B(t)},
\end{align*}
plus remainder cross-terms.

\medskip
\noindent\textbf{Term I.}\enspace
\[
  A(t+\delta)\,\overline{A(t)}
  = \sum_{m,n\le N(t)} \frac{1}{\sqrt{mn}}\,(n/m)^{it}\,m^{-i\delta}.
\]
The diagonal $m=n$ contributes
\begin{equation}\label{eq:DI}
  D_{\mathrm{I}}(t) = \sum_{n\le N(t)} \frac{e^{-i\delta_h(t)\log n}}{n}.
\end{equation}
The off-diagonal terms ($m\ne n$) carry the factor $(n/m)^{it}$ with
$\log(n/m)\ne 0$. By the Montgomery--Vaughan mean-value theorem
\cite{MontgomeryVaughan1974}, since the frequencies $\{\log n : n\le N(t)\}$
are distinct and the window length $T_t \gg N(t)^2$ (as $T_t \asymp \sigma(T)
\gg N(t)^2 \asymp t/(2\pi)$ is false for large $t$; rather $T_t \sim t$ and
$N(t)\sim t^{1/2}$, so $T_t \gg N(t)$), the Ces\`aro contribution of
off-diagonals is
\[
  O\!\left(\frac{N(t)}{T_t}\right)^{1/2}\cdot O(\log T_t)
  = O\!\left(\frac{1}{T_t^{1/2}\log T_t}\right) = o(1).
\]

For the diagonal, from Lemma~\ref{lem:delta},
$\delta_h(t)\log n = h\log n/\vartheta'(t) + O(\log n/(t(\log t)^3))$.
Define the spectral variable $\omega_n(t) \assign -\log n/\vartheta'(t) \in [-1-O(t^{-2}),\,0]$
(using $\log n \le \log N(t) \le \tfrac{1}{2}\log(t/2\pi) = \vartheta'(t) + O(t^{-1})$).
By the Euler--Maclaurin summation formula applied with $f(u) = e^{-iuh}$ and
$u = \log n/\log N(t) \in [0,1]$:
\[
  D_{\mathrm{I}}(t) = \sum_{n\le N(t)}\frac{e^{i\omega_n(t)h(1+O(1/(t\log^2 t)))}}{n}
  = \vartheta'(t)\int_{-1}^0 e^{i\omega h}\,d\omega + O(1),
\]
where the $O(1)$ error from Euler--Maclaurin is $O(\log N(t)/N(t)) = o(1)$ but
contributes $O(1)$ after integration since the Ces\`aro normalisation
$1/\vartheta(T_t)$ times $\int_{t_0}^{T_t} O(1)\,dt = O(T_t)$ gives $O(T_t/\vartheta(T_t))
= O(1/\log T_t) \to 0$. Therefore:
\[
  \frac{1}{\vartheta(T_t)}\int_{t_0}^{T_t} D_{\mathrm{I}}(t)\,dt
  = \frac{1}{\vartheta(T_t)}\int_{t_0}^{T_t}\vartheta'(t)\,dt\cdot\int_{-1}^0 e^{i\omega h}\,d\omega
  + O\!\left(\frac{1}{\log T_t}\right)
  \;\to\; \int_{-1}^0 e^{i\omega h}\,d\omega.
\]

\medskip
\noindent\textbf{Term IV.}\enspace
By the mean value theorem and Lemma~\ref{lem:delta},
\[
  \vartheta(t+\delta) - \vartheta(t) = \delta\,\vartheta'(t) + O(\delta^2\vartheta''(t))
  = h + O\!\left(\frac{1}{t(\log t)^2}\right),
\]
so the phase factor is
\[
  e^{-2i\vartheta(t+\delta)+2i\vartheta(t)} = e^{-2ih}\!\left(1 + O\!\left(\frac{1}{t(\log t)^2}\right)\right).
\]
Term IV is
\[
  e^{-2ih(1+o(1))}\sum_{m,n\le N(t)}\frac{1}{\sqrt{mn}}\,(m/n)^{it}\,m^{i\delta}.
\]
The diagonal $m=n$ contributes
\[
  D_{\mathrm{IV}}(t) = e^{-2ih(1+o(1))}\sum_{n\le N(t)}\frac{e^{i\delta_h(t)\log n}}{n}.
\]
The phase per summand is $e^{i\delta_h(t)\log n} \approx e^{-i\omega_n(t)h}$
(sign reversal from $n^{+it}$ in $B$). Combined with $e^{-2ih}$, the total phase is
$e^{-2ih}\cdot e^{-i\omega_n h} = e^{i(-\omega_n-2)h}$.
Setting $\tilde\omega_n \assign -\omega_n - 2 \in [-2,-1]$ for $\omega_n \in [-1,0]$,
the same Euler--Maclaurin argument gives
\[
  D_{\mathrm{IV}}(t) = \vartheta'(t)\int_{-2}^{-1} e^{i\omega h}\,d\omega + O(1),
\]
and the same Ces\`aro averaging yields
\[
  \frac{1}{\vartheta(T_t)}\int_{t_0}^{T_t} D_{\mathrm{IV}}(t)\,dt
  \;\to\; \int_{-2}^{-1} e^{i\omega h}\,d\omega.
\]

\medskip
\noindent\textbf{Terms II and III.}\enspace
Term II carries the phase $e^{2i\vartheta(t)}$ and Term III carries
$e^{-2i\vartheta(t+\delta)} = e^{-2i\vartheta(t)}e^{-2ih+o(1)}$.
In both cases the phase $e^{\pm 2i\vartheta(t)}$ has instantaneous frequency
$\pm 2\vartheta'(t) \sim \pm\log(t/2\pi) \to \infty$.
The amplitudes are products of Dirichlet polynomials of length $N(t) \asymp t^{1/2}$,
bounded in $L^2$ mean by $O(\vartheta'(t))$ via Montgomery--Vaughan.
By the van der Corput lemma (first derivative test) \cite{Titchmarsh1986},
with phase $\Phi(t) = \pm 2\vartheta(t)$ satisfying $\Phi'(t) = \pm 2\vartheta'(t)
\gg \log t$ and amplitude $F(t) = O(t^{1/2+\varepsilon})$ (from the convexity
bound $|A(t)|, |B(t)| \ll t^{1/4+\varepsilon}$):
\[
  \int_{t_0}^{T_t} F(t)\,e^{\pm 2i\vartheta(t)}\,dt
  \ll \frac{\sup|F(t)|}{\inf|\Phi'(t)|}
  \ll \frac{T_t^{1/2+\varepsilon}}{\log t_0} = O\!\left(\frac{T_t^{1/2+\varepsilon}}{\log T_t}\right).
\]
After Ces\`aro averaging by $1/\vartheta(T_t) \asymp 1/(T_t\log T_t)$:
\[
  \frac{1}{\vartheta(T_t)}\int_{t_0}^{T_t}(\text{Term II or III})\,dt
  = O\!\left(\frac{T_t^{1/2+\varepsilon}}{T_t\log^2 T_t}\right) = o(1).
\]

\medskip
\noindent\textbf{Remainder cross-terms.}\enspace
Every cross-term involving $R(t) = O(t^{-1/4})$ is bounded by
$|R(t)|\cdot(|A(t)|+|B(t)|)$. By Cauchy--Schwarz,
\[
  \int_{t_0}^{T_t}|R(t)|\,|A(t)|\,dt
  \le \left(\int_{t_0}^{T_t}|R(t)|^2\,dt\right)^{1/2}
      \left(\int_{t_0}^{T_t}|A(t)|^2\,dt\right)^{1/2}.
\]
Since $|R(t)| \le C\,t^{-1/4}$, the first factor is $O(T_t^{1/4})$.
By Montgomery--Vaughan, $\int_{t_0}^{T_t}|A(t)|^2\,dt = T_t\sum_{n\le N(t)}n^{-1}
+ O(N(t)^2) = O(T_t\log T_t)$, so the second factor is $O((T_t\log T_t)^{1/2})$.
Thus the cross-term integral is $O(T_t^{3/4}(\log T_t)^{1/2})$, and after
Ces\`aro averaging:
\[
  \frac{O(T_t^{3/4}(\log T_t)^{1/2})}{\vartheta(T_t)}
  = O\!\left(\frac{T_t^{3/4}}{T_t(\log T_t)^{1/2}}\right)
  = O\!\left(\frac{1}{T_t^{1/4}(\log T_t)^{1/2}}\right) = o(1).
\]
The cross-terms involving $\overline{R}$ are handled identically.

\subsection*{Step 4: Assembly}

Combining all terms via Lemma~\ref{lem:jacobian}:
\begin{align*}
  \lim_{T_t\to\infty}\frac{1}{\vartheta(T_t)}\int_{t_0}^{T_t}
  \zeta\!\left(\tfrac{1}{2}+i(t+\delta_h(t))\right)
  \overline{\zeta\!\left(\tfrac{1}{2}+it\right)}\,dt
  &= \int_{-1}^0 e^{i\omega h}\,d\omega
   + \int_{-2}^{-1} e^{i\omega h}\,d\omega \\
  &= \int_{-2}^0 e^{i\omega h}\,d\omega
   = \frac{2e^{-ih}\sin h}{h}.
\end{align*}
Hence $\lim_{T\to\infty}\mathcal{I}(T,h) = 2e^{-ih}\sin h/h$, and by
\eqref{eq:split} and Lemma~\ref{lem:E-vanishes}:
\[
  R(h) = \tfrac{1}{2}\cdot\frac{2e^{-ih}\sin h}{h} = \frac{e^{-ih}\sin h}{h}.
\]

\noindent\textbf{Consistency check at $h=0$.}\enspace
$R(0) = \lim_{h\to 0}(e^{-ih}\sin h)/h = 1$, consistent with unit
Ces\`aro mean square of $X$. \qed

\begin{thebibliography}{9}

\bibitem{Ingham1928}
A.\,E.~Ingham,
\textit{Mean-value theorems in the theory of the Riemann zeta-function},
Proc.\ London Math.\ Soc.\ (2) \textbf{27} (1928), 273--300.

\bibitem{Titchmarsh1986}
E.\,C.~Titchmarsh,
\textit{The Theory of the Riemann Zeta-Function},
2nd ed.\ (revised by D.\,R.~Heath-Brown),
Oxford University Press, 1986.

\bibitem{MontgomeryVaughan1974}
H.\,L.~Montgomery and R.\,C.~Vaughan,
\textit{Hilbert's inequality},
J.\ London Math.\ Soc.\ (2) \textbf{8} (1974), 73--82.

\bibitem{Paris2020}
R.\,B.~Paris,
\textit{Refined asymptotics of the Riemann--Siegel theta function},
arXiv:2004.00926, 2020.

\end{thebibliography}

\end{document}
